\section*{How I Found Out About the Mentorship Program}

I discovered the OpenSSF LFX Mentorship Program directly through the LFDT website, where I actively search for structured open-source opportunities. The program's focus on impactful projects within supply chain security, like the Hyperledger Fabric-X and Fablo integration, immediately captured my attention. I find these programs crucial for gaining hands-on experience in real-world infrastructure rather than isolated academic exercises.

% Divider line after question 1
\begin{center}
    \textcolor{secondary}{\rule{0.7\textwidth}{0.3pt}}
\end{center}

\section*{Why I Am Interested in This Program}

My interest stems from a long-standing engagement with open-source and decentralized technologies, seeking to apply my skills to impactful hands-on development. The Fablo project's role in lowering the barrier to entry for Hyperledger Fabric-X deeply resonates with me, highlighting its community importance. This ambitious integration offers a significant learning curve, promising substantial growth in my understanding of complex systems and the satisfaction of contributing to a tool that benefits many.

% Divider line after question 2
\begin{center}
    \textcolor{secondary}{\rule{0.7\textwidth}{0.3pt}}
\end{center}

\section*{Experience and Knowledge/Skills Applicable to This Program}

My two years of open-source contributions, marked by over 800 GitHub contributions, have provided extensive experience in distributed systems, Docker-based tooling, TypeScript/Node.js, and CI/CD workflows, all directly relevant to this project. I have leveraged Docker and GitHub Actions for containerized applications and multi-service orchestration, crucial for managing Fabric-X networks via Docker Compose and Fablo. Specifically, my prior study of Fablo’s architecture and the Fabric-X token sample, including its xdev topology, core.yaml, routing tables, and fxconfig lifecycle, culminated in a functional MVP implementing `global.platform: "fabricx"` generator branching and network lifecycle commands. This focused work positions me exceptionally well to drive the Fabric-X integration into Fablo.

% Divider line after question 3
\begin{center}
    \textcolor{secondary}{\rule{0.7\textwidth}{0.3pt}}
\end{center}

\section*{What I Hope to Get Out of This Mentorship Experience}

Through this mentorship, I aim to significantly enhance my proficiency with complex codebase architectures and deepen my understanding of both Fablo and Fabric-X. I particularly look forward to gaining advanced confidence in Docker-based environments and microservice-oriented systems. Learning directly from the mentors' real-world experience and guidance is paramount for my development as an effective open-source contributor, enabling me to make lasting, meaningful contributions to the community and the project's future evolution.

% Divider line after question 4
\begin{center}
    \textcolor{secondary}{\rule{0.7\textwidth}{0.3pt}}
\end{center}

\vspace{0.5cm}

% Sincerity section (left-aligned, bigger name)
\begin{flushleft}
    \textbf{\color{primary}Sincerely,} \\
    \vspace{0.2cm}
    {\color{primary}\Large \textbf{Harshit Kandpal}} \\
    \vspace{0.1cm}
    {\color{secondary}B.Tech Computer Science, IIT Bhilai (2024–2028)}
\end{flushleft}